\startvolume{Paper VII}{Classical recoupling packets}

\lesson{Paper VII / 01}
{A ladder begins with a positive shape}
{vii_01_ladder}
{K_0|n\rangle=(n+\kappa)|n\rangle,\qquad K_+|n\rangle=\sqrt{(n+1)(n+2\kappa)}\,|n+1\rangle}
{Read the left panel vertically: an occupation label selects a weight eigenvalue. Changing the positive shape shifts that entire ladder. The right panel reads the same labels as raising amplitudes. Those coefficients determine how neighboring degrees are joined, rather than merely naming them. The curves are computed from the defining representation, with no fitted model. These two pieces of data are the input to every tensor construction in this volume.}
{Set $\kappa=1/2$ and start at degree $n=2$. Its weight is $2.5$, while its raising coefficient is $\sqrt{3\cdot3}=3$. Thus $K_+|2\rangle=3|3\rangle$ in the algebraic ladder. This coefficient is an operator amplitude, not a probability. A normalized physical transition would require its own state and dynamics. The example explains why later Casimir matrices retain products of ladder amplitudes.}
{Keep the degree label, weight eigenvalue, and transition amplitude as three distinct quantities.}
{VII, Chapter 1, Definition 1.1, \emph{Positive shape packet}, source lines 219--243. Direct representation formulas; positive $\kappa$, nonnegative integer degree, finite-support vectors.}

\lesson{Paper VII / 02}
{Project a large excitation simplex}
{vii_02_simplex}
{\dim H_N^{(3)}=\binom{N+2}{2},\qquad \dim M_N^{(3)}=N+1}
{Each point on the left is one occupation triple with the same total excitation. The three digits give its exact address, and color records the third occupation. The plot contains fifteen points because weak compositions count the ambient basis. On the right, compare that full count with the dimension of the new lowest-weight kernel. The smaller curve counts an exact subspace. It does not mean that a few occupation states were discarded by sampling.}
{At $N=4$, the full space has $\binom62=15$ basis states. The lowering map targets degree three, whose full dimension is ten. Its full row rank therefore leaves a five-dimensional kernel. A subsequent recoupling matrix needs $5\times5$ entries, even though its defining projector was constructed in fifteen dimensions. At $N=0$, both spaces have dimension one, so the reduction also respects the initial degree.}
{Compression comes from an operator kernel, with its dimension proved before any numerical atlas is chosen.}
{VII, Chapter 2, Proposition 2.1 and Theorem 2.2. The explicit full-row-rank proof used here is VIII, Theorem 3.1, specialized at $q=1$.}

\lesson{Paper VII / 03}
{See the lowest-weight plane}
{vii_03_lowest_plane}
{M_N=\ker(K_-^{[3]}|_{H_N}),\qquad N=1:\quad x_{100}+x_{010}+x_{001}=0}
{The three-dimensional view is coefficient space. Its axes are amplitudes of three occupation basis states, rather than spatial coordinates of particles. The shaded plane is the exact lowering kernel for three half-shapes at one excitation. Two arrows show an orthonormal atlas inside that plane. The right panel advances one degree and shows the rectangular lowering matrix. The next kernel is again selected by an equation, not by the apparent proximity of points in a picture.}
{At degree one, $(1,1,-2)/\sqrt6$ lies in the plane because its coefficients sum to zero. The vector $(1,-1,0)/\sqrt2$ also lies there, and their inner product vanishes. At degree two, the displayed $3\times6$ lowering map has a three-dimensional nullspace. The generated frame satisfies its lowering equation to floating-point roundoff. This residual checks this frame; the rank argument proves the general dimension.}
{A basis is a choice inside the kernel; the kernel is the mathematical object being represented.}
{VII, Chapter 2, \emph{Fixed excitation and the lowest-weight projector}, lines 259--343; Chapter 5, the one-excitation construction. Newly computed plane and matrix.}

\lesson{Paper VII / 04}
{Two observables select two atlases}
{vii_04_pair_atlases}
{\lambda_j^{ab}=(\kappa_a+\kappa_b+j)(\kappa_a+\kappa_b+j-1),\qquad \mathcal R=L^{\mathsf T}R}
{The left curves are the spectra of two intermediate Casimirs. Their labels record how two neighboring shape factors are coupled. The right panel is the overlap of their complete eigenbases in the same lowest-weight space. A spectrum gives the addresses of an atlas, while the overlap tells how coordinates move between atlases. Eigenvalues alone do not determine that movement. The signed entries retain the phase convention of the selected real eigenvectors.}
{Use shapes $(0.7,1.1,0.9)$ and $N=4$. The left pair begins at $1.8(0.8)=1.44$, and the right pair begins at $2(1)=2$. Both spectra contain five entries because the multiplicity dimension is five. The program restricts each Casimir, orders its eigenvalues, and computes the displayed overlap. It separately compares those eigenvalues with the formula and checks orthogonality. These are complementary numerical checks of the chosen example.}
{The shared space supports reconstruction; the different intermediate observables support different coordinate descriptions.}
{VII, Chapter 3, Proposition 3.1; Chapter 4, Definition 4.1 and Theorem 4.2. Pair-spectrum identification is a classical representation-theoretic input.}

\lesson{Paper VII / 05}
{The middle shape controls a beam splitter}
{vii_05_shape_beam}
{\mathcal R_1=\begin{pmatrix}c&s\\s&-c\end{pmatrix},\qquad c^2=\frac{\kappa_1\kappa_3}{(\kappa_1+\kappa_2)(\kappa_2+\kappa_3)},\quad s^2=1-c^2}
{Hold the outer shapes fixed and move along the middle-shape axis. The first curve is the retained-channel amplitude, and the second is the transferred-channel amplitude in the declared real gauge. Their squares always sum to one. The right panel freezes the curve at one shape triple and displays the full signed transform. Its negative entry is needed for orthogonality. Reading the same numbers as unsigned probabilities would erase the operation that makes exact synthesis possible.}
{For $(\kappa_1,\kappa_2,\kappa_3)=(0.7,1.1,0.9)$, the denominator is $1.8\cdot2=3.6$. Hence $c^2=0.175$ and $s^2=0.825$, giving $c\approx0.41833$ and $s\approx0.90830$. The independent Casimir construction reproduces this matrix. Applying the overlap and then its transpose recovers every two-channel amplitude vector. Squaring entries instead produces a measured channel with a different role.}
{The exact low-degree formula links a continuous shape choice to a finite reversible coordinate transform.}
{VII, Chapter 5, \emph{One excitation as a shape beam splitter}, source lines 452--510. Exact formula checked against the newly computed pair-Casimir overlap.}

\lesson{Paper VII / 06}
{An exact three-channel transform}
{vii_06_exact_three}
{\mathcal R_2=\begin{pmatrix}
1/3&1/\sqrt3&\sqrt5/3\\
1/\sqrt3&1/2&-\sqrt{15}/6\\
\sqrt5/3&-\sqrt{15}/6&1/6
\end{pmatrix},\qquad P_{jm}=|\mathcal R_2(j,m)|^2}
{Read the left matrix as amplitudes and the right matrix as probabilities after measurement. The two negative amplitudes participate in cancellations, while the probability matrix contains only nonnegative entries. Rows and columns of the amplitude matrix are orthonormal. Rows and columns of the probability matrix sum to one. Both panels describe the same exact half-shape example, but they support different operations. The displayed decimals are visual labels for the exact radicals in the central formula.}
{The first amplitude row has squared norm $1/9+1/3+5/9=1$. Its probability row is therefore $(1/9,1/3,5/9)$. For the first two amplitude rows, positive and negative products cancel, giving inner product zero. The generated Casimir overlap agrees with the complete radical matrix. Checking just the first probability row would confirm one normalization, while the full orthogonality check confirms reversible reconstruction across all three channels.}
{Keep the exact signed matrix available even when the application eventually requests only measured probabilities.}
{VII, Chapter 6, \emph{An exact three-channel example}, source lines 512--559, equal shapes $1/2$ and $N=2$.}

\lesson{Paper VII / 07}
{Phase changes a measured output}
{vii_07_phase_measurement}
{b_m=\sum_j\mathcal R_{jm}a_j,\qquad |b_m|^2\ \text{and}\ \sum_j|\mathcal R_{jm}|^2|a_j|^2\ \text{need not agree}}
{The left bars are the two signed contributions reaching each target channel. The black points are their coherent sums. The right panel compares measuring those sums with dephasing the input first and then using the squared-amplitude channel. Both output laws are normalized, yet they differ because only one retains the cross terms between routes. This is a concrete reason to separate a reversible coordinate transform from a stochastic measurement kernel in the packet record.}
{Start in the exact three-channel example with $a=(1,-1,0)/\sqrt2$. At target label zero, the coherent amplitude is $(1/3-1/\sqrt3)/\sqrt2$. Its squared magnitude is approximately $0.02977$. If the input is dephased first, the corresponding probability is $(1/9+1/3)/2=2/9$. The probability difference comes from interference. It is not a failure of normalization, and the signed matrix itself still has its exact inverse.}
{Amplitude addition occurs before squaring in coherent synthesis; a dephased channel has already lost those cross terms.}
{VII, Chapters 9 and 11, measurement interpretation and Boundary 11.1. Newly worked superposition using the exact Chapter 6 matrix.}

\lesson{Paper VII / 08}
{A positive recurrence needs its own gauge}
{vii_08_recurrence_gauge}
{J=\mathcal R\,\operatorname{diag}(\lambda_m)\mathcal R^{\mathsf T},\qquad S=\operatorname{diag}(s_j):\ (SJS)_{j,j+1}=|J_{j,j+1}|}
{Both panels represent the same opposing Casimir. In the left panel, the overlap convention makes its first row and column positive. The resulting adjacent Casimir entries can be negative. The right panel changes the signs of left-atlas vectors so every adjacent entry is positive. The zero pattern and eigenvalues remain unchanged. This distinction matters when a scalar three-term recurrence is presented as a Jacobi recurrence with positive neighboring coefficients.}
{For three half-shapes at $N=1$, the printed overlap has $J_{01}=-\sqrt3/2$. It cannot simultaneously have the same signed coordinates and a positive adjacent coefficient. Set $s_0=1$, then choose $s_{j+1}=s_j\operatorname{sign}(J_{j,j+1})$. The rephased overlap is $S\mathcal R$, and its recurrence uses $SJS$. The displayed seven-channel example checks the tridiagonal zero pattern and the rephased matrix eigen-equation separately.}
{A gauge changes coordinates; state it explicitly before assigning signs to recurrence coefficients.}
{VII, Chapter 7, Theorem 7.1, read with the corrected row-sign convention. Irreducible tridiagonality is a classical Racah input; the displayed finite matrices are newly checked.}

\lesson{Paper VII / 09}
{A fixed measured channel spreads a law}
{vii_09_entropy}
{P_{jm}=|\mathcal R_{jm}|^2,\qquad H(pP)\geq H(p),\qquad H(p)=-\sum_jp_j\log p_j}
{The bars compare one input law with its output under one fixed squared-overlap matrix. The right panel sorts both distributions before adding their largest entries. The output curve lies below the input curve until both reach total mass one. That is the finite concentration order behind the entropy statement. The comparison is between input and output at a fixed kernel. It does not compare unrelated kernels or promise that entropy rises when a shape parameter is changed.}
{At $N=5$, use shapes $(0.7,1.1,0.9)$ and the input $(0.7,0.1,0.05,0.05,0.05,0.05)$. The code computes its six-channel measured output directly. Each row and column of the kernel sums to one, and the output entropy exceeds the input entropy. Concavity applied with the column weights proves the inequality for every probability law. No random sampling is needed for that theorem.}
{Bistochasticity gives a universal fixed-channel entropy law, while particular mixing rates remain properties of the chosen kernel.}
{VII, Chapter 9, Theorem 9.1; Chapter 10, Theorem 10.1. Natural logarithms give entropy in nats; zero terms use $0\log0=0$.}

\lesson{Paper VII / 10}
{Tree coordinates have joint addresses}
{vii_10_tree_addresses}
{\dim M_N^{(4)}=\binom{N+2}{2},\qquad (((12)3)4):\quad 0\leq j_{12}\leq j_{123}\leq N}
{Each vertex on the left names a binary bracketing of four factors. The two block labels beneath it name the commuting intermediate Casimirs used by that tree. The right panel shows their joint integer addresses for the first tree. These addresses fill a triangular lattice rather than an unconstrained square. An atlas is complete only when all compatible addresses occur in the same final lowest-weight space. The graph explains which coordinate transforms connect neighboring bracketings.}
{At $N=4$, there are fifteen addresses. With shapes $(0.55,0.8,1.05,1.3)$, the address $(j_{12},j_{123})=(1,3)$ gives intermediate shapes $2.35$ and $5.4$. Their ordinary Casimir eigenvalues are $3.1725$ and $23.76$. The numerical joint eigenbasis recovers the full triangular address set and diagonalizes both operators. Dimension, address completeness, and diagonalization each test a different part of this construction.}
{A joint spectrum supplies a finite address system for a tree, rather than just a list of separate eigenvalues.}
{VII, Chapter 12, \emph{Five complete coupling trees} and Proposition 12.1; plotted addresses are extracted from newly computed joint eigenbases.}

\lesson{Paper VII / 11}
{What a pentagon check actually certifies}
{vii_11_loop_scope}
{U_{r,r+1}=B_r^{\mathsf T}B_{r+1},\qquad U_{01}U_{12}U_{23}U_{34}U_{40}=I}
{The left bars show that three constructions all close their common-space overlap loop. The first uses the intended Casimir atlases. The second changes coordinates consistently at every vertex. The third uses a common subspace that fails the intended lowering equation. The right bars reveal that difference. The loop theorem remains exact, but its telescoping proof cannot independently distinguish the intended projector from every consistently shared alternative. The plotted log scale displays roundoff and a large failure together.}
{For four factors at $N=2$, the intended packet dimension is six inside a ten-dimensional excitation space. Five complete orthonormal bases of a common six-dimensional subspace always telescope. This remains true after consistent orthogonal coordinate changes. Choosing six ambient coordinate vectors also produces a closing loop, yet their lowering residual exceeds one. Thus the loop check passes while the separate kernel check rejects the wrong subspace.}
{Pair overlap closure with kernel, dimension, spectrum, and eigen-equation checks when validating an associator implementation.}
{VII, Chapter 13, Theorem 13.1, with its implementation-test interpretation corrected. Newly constructed deterministic negative control, coordinate-change seed 73011.}

\lesson{Paper VII / 12}
{Compile in the reduced dimension}
{vii_12_compiler}
{\mathcal R=G_1^{\mathsf T}\cdots G_L^{\mathsf T}D,\qquad L\leq d(d-1)/2,\qquad D_{jj}\in\{-1,1\}}
{The left panel is the signed diagonal remaining after row rotations eliminate entries below the overlap's diagonal. Reversing those rotations reconstructs the original transform. The right panel compares generated rotation counts with the triangular worst-case bound as the packet dimension grows. Each rotation touches only two coordinate slots. This compiler uses the multiplicity dimension, so its index set is smaller than the original tensor excitation space that defined the projector.}
{At $N=5$ for three factors, the packet dimension is $d=6$ and the bound is fifteen rotations. The program records each Givens matrix, reduces the overlap to diagonal signs, and reconstructs it in reverse order. The residual is checked against the original overlap. This gives a finite algebraic circuit for the example. It does not assert nearest-neighbor hardware locality, minimum circuit depth, or an optimal gate count for a particular device.}
{Reduction defines the working coordinates; rotation factorization supplies an exact reconstruction procedure in those coordinates.}
{VII, Chapter 15, Theorem 15.1, \emph{Givens packet compilation}; Chapter 16, \emph{Projector-resolved Racah transform}. Direct finite linear algebra and new bounded checks.}

\startvolume{Paper VIII}{q-deformation and scalar bridges}

\lesson{Paper VIII / 01}
{Replace a linear number by a q-number}
{viii_01_q_numbers}
{[x]_q=\frac{q^x-q^{-x}}{q-q^{-1}}=\frac{\sinh(x\theta)}{\sinh\theta},\qquad q=e^{-\theta},\quad [x]_1=x}
{The left panel shows the symmetric q-number at three positive deformation choices. The ordinary linear coordinate is recovered at one. The right panel divides by the positive input, making the relative deformation visible. This convention is symmetric under replacing q by its reciprocal, even though the ordered coproduct later retains directional q-weights. A scalar convention must be fixed before formulas from different quantum-algebra sources are compared.}
{At $q=1/2$, $[2]_q=2.5$ and $[3]_q=5.25$, while $[1]_q=1$. Thus the half-shape ladder's raising amplitude at degree one becomes $\sqrt{[2]_q[2]_q}=2.5$. At $q=1$, that coefficient is two. The stable implementation uses the hyperbolic expression near one and the exact continuous value at one. Positive x always gives a positive q-number in the declared real range.}
{Deforming the scalar ladder coefficients preserves the positive shape construction while changing its degree-dependent geometry.}
{VIII, Chapter 1, Definitions 1.1--1.2 and Proposition 1.3. Real range $0<q\leq1$; q=1 is defined by continuity.}

\lesson{Paper VIII / 02}
{The coproduct routes lowering events}
{viii_02_coproduct}
{\Delta K_-=K_-\otimes q^{K_0}+q^{-K_0}\otimes K_-,\qquad (\Delta\otimes I)\Delta=(I\otimes\Delta)\Delta}
{Each curve on the left corresponds to lowering a different occupied factor at total degree one. The labels follow the declared composition order. Their unequal weights come from the q-powers assigned to factors before and after the lowering event. On the right, an explicit three-term expansion agrees with the iterated matrix implementation. This tests coassociative routing in a small transparent case. Coassociativity concerns bracketings of the same ordered factors, rather than freely permuting those factors.}
{At shapes $(0.7,1.1,0.9)$ with $q=0.73$, the $|001\rangle$ route has coefficient $q^{-(0.7+1.1)}\sqrt{[1.8]_q}$. The route from $|010\rangle$ has coefficient $q^{-0.7+0.9}\sqrt{[2.2]_q}$, and the $|100\rangle$ route has $q^{1.1+0.9}\sqrt{[1.4]_q}$. Both coproduct bracketings produce exactly these three terms. Their values also match the freshly built lowering row.}
{Preserve the tensor order and q-number convention when constructing block observables or comparing formulas.}
{VIII, Chapter 2, Theorem 2.1, \emph{Coassociative packet routing}, source lines 293--327. Direct expansion and new degree-one matrix comparison.}

\lesson{Paper VIII / 03}
{A q-Hahn seed keeps signs and weights}
{viii_03_lowest_seed}
{\frac{c_{n+1}}{c_n}=-q^{-(N-1+\kappa_1+\kappa_2)}
\sqrt{\frac{[N-n]_q[2\kappa_2+N-n-1]_q}{[n+1]_q[2\kappa_1+n]_q}}}
{The left panel retains the alternating phases of the one-dimensional lowest-weight seed. The right panel squares those same coefficients and exposes the q-dependent weight tilt. Both come from a solved lowering equation. Changing q alters the distribution without changing the six available occupation labels. The normalized vectors have total probability one. Identifying general coupled coefficients with scalar q-Hahn polynomials is an additional classical representation-theoretic input.}
{Take $N=5$ and shapes $(0.7,1.1)$. Starting with $c_0=1$, apply the displayed ratio five times, then normalize. At $q=0.73$, the resulting vector is annihilated by the two-factor lowering matrix. It also agrees up to its overall sign with a separately computed nullspace vector. The ratio derives the seed analytically; the two numerical checks confirm its implementation at this input. The plotted squared weights are not inferred from a fitted probability model.}
{The signed seed supports synthesis, and its square supports a weight description; neither representation should replace the other silently.}
{VIII, Chapter 4, Theorem 4.1, \emph{Explicit q-Hahn lowest packet}, source lines 407--474. New recurrence and nullspace comparisons at three q choices.}

\lesson{Paper VIII / 04}
{Deform both a spectrum and an atlas}
{viii_04_deformed_atlases}
{\lambda_j^{ab}(q)=[\kappa_a+\kappa_b+j]_q[\kappa_a+\kappa_b+j-1]_q,\qquad \mathcal R_N(q)=L_N(q)^{\mathsf T}R_N(q)}
{The left curves follow one pair's intermediate eigenvalues as q changes. On the right, the complete overlap at one q value records how the left and right Casimir atlases mix. The packet dimension remains seven throughout this example. The changed colors describe coordinate geometry on a fixed-size object, rather than a change in the number of channels. Spectral agreement and overlap orthogonality are checked separately.}
{For $N=6$, shapes $(0.7,1.1,0.9)$ and $q=0.73$, build the full degree-six lowering map, compute its seven-dimensional kernel, and restrict both pair Casimirs. Their ordered eigenvectors define the overlap. The code compares the left eigenvalues with the displayed q-number formula and checks $\mathcal R^{\mathsf T}\mathcal R=I$. This is a finite matrix construction. The exact scalar q-Racah parameter and weight map is kept as a separate comparison obligation.}
{The matrix module is usable before an explicit scalar normalization formula is supplied, provided its evidence is described at that level.}
{VIII, Chapter 5, Theorem 5.1 and \emph{Two intermediate q-Casimir atlases}. Classical pair-spectrum/q-Racah inputs; the plotted matrix is newly computed.}

\lesson{Paper VIII / 05}
{An exact q-coordinate reaches a swap}
{viii_05_q_beam}
{\mathcal R_q=\begin{pmatrix}c_q&s_q\\s_q&-c_q\end{pmatrix},\qquad
c_q=\frac{q}{1+q^2},\quad s_q=\frac{\sqrt{1+q^2+q^4}}{1+q^2}}
{Read the amplitude curves from small q toward one. The diagonal amplitude rises from zero to one half, while the transferred amplitude approaches $\sqrt3/2$. The right panel displays the complete signed transform at the midpoint q value. Its limit at zero is a routing swap. This is an exact limit for this normalized two-channel model. The unnormalized q-number operators themselves do not have a uniform finite limit at zero.}
{At $q=0.5$, $c_q=0.4$ and $s_q=\sqrt{0.84}\approx0.916515$. Their squares sum to one, and the lower-right sign makes the two columns orthogonal. The numerical pair-Casimir construction matches this matrix at four q values. At $q=1$, it recovers the classical half-shape beam splitter. At a very small positive q, the closed expression approaches the swap matrix directly, avoiding an ill-conditioned large-degree operator construction.}
{State the normalization and tested model before promoting a visually sparse q endpoint into a general routing theorem.}
{VIII, Chapter 6, Theorem 6.1, \emph{Closed two-channel q-Racah transform}, half-shapes and $N=1$. Routing concerns this closed formula only.}

\lesson{Paper VIII / 06}
{Two-channel mixing has a monotone gap}
{viii_06_entropy_gap}
{p(q)=\frac{q^2}{(1+q^2)^2},\qquad \lambda_2=2p-1,\qquad
\operatorname{gap}_{\rm abs}=1-|\lambda_2|=2p}
{The left curve is the binary entropy of one measured row. The right panel separates the absolute mixing gap from the magnitude of the other eigenvalue. Small q is close to a periodic swap, whose second eigenvalue has magnitude one. Moving toward q=1 reduces that magnitude and increases the absolute gap. The word absolute matters: a negative second eigenvalue still produces an oscillating mode, and its distance from one alone is not the mixing rate used here.}
{At $q=0.5$, the measured matrix has rows $(0.16,0.84)$ and $(0.84,0.16)$. Its second eigenvalue is $-0.68$, its absolute gap is $0.32$, and its row entropy is approximately $0.43967$ nats. The smaller row probability increases only to $1/4$ over the declared interval. Binary entropy is increasing there, which proves the model's q-monotonicity without relying on the plotted grid.}
{The exact one-excitation model has a monotone deformation law; higher multiplicity requires its own analysis.}
{VIII, Chapter 7, Theorem 7.2, \emph{Monotone one-excitation mixing}. Equal half-shapes, $N=1$, $0<q\leq1$, absolute spectral gap.}

\lesson{Paper VIII / 07}
{Higher-level row entropy can decrease}
{viii_07_nonmonotone}
{H_1(q)=-\sum_m|\mathcal R_3(q)_{1m}|^2\log|\mathcal R_3(q)_{1m}|^2,\qquad
H_1(q_1)-H_1(q_0)\approx-0.007784}
{The left panel follows the four probabilities of the second measured row, using zero-based row index one. The right panel computes their entropy and highlights two nearby q values. The curve decreases between them. Every individual kernel remains bistochastic, so each still obeys the input-to-output entropy inequality. The decrease compares two different kernels' row laws. It therefore refutes universal entropy monotonicity in q without contradicting the fixed-channel theorem.}
{For $N=3$ and three half-shapes, take $q_0=307/395\approx0.777215$ and $q_1=1557/1975\approx0.788354$. The entropies are approximately $1.244882$ and $1.237098$ nats. Their difference is negative by about $0.007784$. The guide recomputes both matrices in double precision from the q-coproduct. A separate full-space implementation records a decreasing interval from $q=0.78$ to $0.80$. These are bounded numerical witnesses, not interval proof certificates or a classification of all deformation behavior.}
{A negative control defines the theorem boundary as precisely as a successful example defines its mechanism.}
{VIII, Chapter 7, Boundary 7.3, \emph{No universal entropy monotonicity in q}. New endpoint witness in vii\_viii\_audit.json and separately saved independent full-space evidence.}

\lesson{Paper VIII / 08}
{Differentiate a complete coordinate transform}
{viii_08_tangent}
{A(\theta)=\mathcal R(\theta)^{\mathsf T}\partial_\theta\mathcal R(\theta),\qquad
A^{\mathsf T}=-A,\quad \partial_\theta\mathcal R=\mathcal R A}
{The left matrix is the skew tangent generator for the seven-channel transform at $N=6$. Opposite entries have opposite signs, and diagonal entries vanish. The right curves measure squared speed for two excitation levels in the same $\theta=-\log q$ coordinate. This quantity concerns sensitivity of the complete atlas overlap, rather than entropy of its measured rows. Its size depends on parameterization, while constant orthogonal changes of atlas preserve the Frobenius coefficient.}
{At $\theta=0.35$, compute the overlaps at $\theta\pm2\times10^{-5}$, align their real sign conventions, and form a central difference. Before projecting to a skew matrix, the code checks the raw generator's symmetric residual. It also checks the tangent equation afterward. This avoids treating an explicitly antisymmetrized output as independent evidence of skewness. Analytically, differentiating $\mathcal R^{\mathsf T}\mathcal R=I$ proves the exact skew law in a smooth gauge.}
{A deformation tangent is meaningful only after the atlas conventions are made continuous along the chosen coordinate.}
{VIII, Chapter 8, Theorem 8.1, \emph{Skew deformation generator}. New finite-difference diagnostics; derivative step and raw residual recorded in the audit.}

\lesson{Paper VIII / 09}
{Approach the classical atlas carefully}
{viii_09_classical_limit}
{[x]_q=x+O((\log q)^2),\qquad \mathcal R_N(q)\longrightarrow\mathcal R_N(1)
\quad\text{in a matched sign gauge}}
{Both curves on the left tend toward zero error as q approaches one. One compares complete recoupling overlaps, and the other compares normalized two-body seed vectors. Their rates differ here. The right panel retains the signed overlap correction at $q=0.99$. The q-number's even hyperbolic expansion does not by itself establish the same rate for every seed or coordinate choice, because the coproduct also carries directional q-powers.}
{For $N=6$ and shapes $(0.7,1.1,0.9)$, compute the q=1 overlap as the reference. Align row and column signs before taking entry differences at nearby q. The seed comparison uses the two-body degree-five vector for shapes $(0.7,1.1)$. Both finite comparisons converge over the displayed range. They verify these numerical examples and conventions. The continuity argument is supplied by the finite matrix construction and its nondegenerate spectral data.}
{Match atlases before measuring a limit, and distinguish an operator coefficient expansion from an entire reconstructed object's convergence rate.}
{VIII, Chapter 9, \emph{Classical and routing limits}, source lines 751--763. New q-grid comparisons at fixed positive shapes; no universal rate asserted.}

\lesson{Paper VIII / 10}
{q-trees need more than loop closure}
{viii_10_q_trees}
{U_{r,r+1}(q)=B_r(q)^{\mathsf T}B_{r+1}(q),\qquad
U_{01}U_{12}U_{23}U_{34}U_{40}=I\quad\text{(indices modulo 5)}}
{The five vertices still represent the same ordered planar bracketings after deformation. Each now uses q-Casimirs on contiguous blocks. The right panel plots three separate diagnostics: common-space loop closure, annihilation by total lowering, and joint diagonalization. Their residuals remain small in the displayed finite tests. They answer different questions. Combining them gives more information than calling the closing overlap loop alone a certificate of every label, phase, and projector convention.}
{At four factors and $N=3$, the multiplicity dimension is ten. For q values $0.65$, $0.75$, $0.85$, and one, the code constructs every joint atlas and composes its five overlaps. It separately tests commuting intermediate operators and the displayed kernel and diagonalization residuals. Consistent coordinate changes would still leave the loop equal to identity. A future stronger local-associator check would compare independently constructed q-Racah edge blocks with these global overlaps before composing them.}
{Keep structural operator checks alongside loop checks when the same data construction could make closure automatic.}
{VIII, Chapter 10, Theorem 10.1, \emph{Finite q-pentagon certificate}, with corrected interpretation. Four-factor inputs and residuals recorded separately.}

\lesson{Paper VIII / 11}
{A scalar table is a separate comparison}
{viii_11_scalar_bridge}
{\begin{gathered}
R_n(\mu_{\gamma,\delta}(y);\alpha,\beta,\gamma,\delta\mid q)
=R_y(\mu_{\alpha,\beta}(n);\gamma,\delta,\alpha,\beta\mid q)\\
\mu_{u,v}(t)=q^{-t}+uvq^{t+1}
\end{gathered}}
{The left panel is generated from the terminating scalar q-hypergeometric expression, with independent scalar parameters. Exchanging degree and coordinate together with the parameter pairs gives the recorded duality agreement. This table has not been identified entry by entry with the shape-Casimir overlap. The right panel gives a finite-support counterexample: the two lattice points coincide, and two evaluated polynomial rows coincide. A denominator truncation condition alone therefore cannot guarantee a positive orthogonal family.}
{Set $N=1$, $q=1/2$, $\alpha=\beta=1/2$, $\gamma=4$, and $\delta=1$. Although $\gamma q=q^{-N}$, both lattice values are three. The finite formula gives $R_0=1$ and $R_1(\mu)=6-(5/3)\mu$, so both row vectors are $(1,1)$. Their inner product under positive weights is $w_0+w_1>0$. This exact rational witness separates termination, nondegeneracy, and positivity. The complete shape-to-scalar map still needs base, parameters, weights, norms, and phases.}
{Scalar identities can be verified independently while a matrix-to-polynomial normalization comparison remains explicit unfinished work.}
{VIII, Chapter 13, Theorem 13.1; Chapter 19, first open direction. Scalar sums require $0<q<1$ and nonzero finite denominator factors; $q=1$ is a separate limit. The collapsed-lattice witness uses exact rational arithmetic.}

\lesson{Paper VIII / 12}
{Check a shifted Askey--Wilson operator}
{viii_12_askey_wilson}
{\mathcal L\mathcal A_n=(q^{-n}+abcdq^{n-1})\mathcal A_n,\qquad
\mathcal L=D_q+(1+abcd/q)I}
{The left curves are normalized terminating Askey--Wilson polynomials on the continuous unit-circle coordinate. The right panel substitutes six degrees into the q-difference equation at thirteen nonsingular points. The residual scale shows numerical agreement. The eigenvalue uses the shifted operator: $D_q$ is the two-difference part, while the scalar identity term has been added. Omitting that shift would change the eigenvalue convention and make the comparison inconsistent.}
{Take $q=0.7$ and $(a,b,c,d)=(0.2,0.3,-0.25,0.4)$. At degree two, the shifted eigenvalue is $q^{-2}+abcdq\approx2.036616$. Evaluate the finite polynomial at $z$, $qz$, and $z/q$, then apply the declared operator coefficients away from their poles. The program compares this result with the eigenvalue times the unshifted polynomial value. This checks a classical scalar identity. It does not construct a continuous-series unitary packet transform or its spectral measure.}
{A continuous variable in a scalar polynomial is a bridge to further analysis, not an already completed continuous-series transform.}
{VIII, Chapter 14, Theorem 14.1, \emph{Askey-Wilson eigenfunction law}. Scalar sums require $0<q<1$, nonzero finite denominator factors, and nonsingular operator points; $q=1$ is a separate limit. Classical identity, newly evaluated on the declared grid.}
